\documentclass[text03.tex]{subfiles}
\begin{document}
\chapter{温度や\ruby{湿度}{しつ|ど}、明るさを\ruby{測}{はか}ってみよう}
\section{はじめに}
\subsection{この章で学ぶこと}
この章では以下のことを学びます。







\begin{itemize}
    \item ターミナルで基本的なコマンドを使う
    \item ファイルとディレクトリの考え方を知る
    \item センサーボードをプログラムで動かす
    \item 温度、湿度、気圧、明るさを測る・LEDやボタンをプログラムで操作する
\end{itemize}

この章の前半では、コンピュータをコマンドで動かす方法を学びます。
マウスではなく、コマンドでコンピュータを\ruby{操作}{そう|さ}することは、とても重要です。
コマンドを使うと、コンピュータがするべき操作を文字で記録したり、
\ruby{離}{はな}れたコンピュータを\ruby{簡単}{かん|たん}に操作できるようになります。

次に、センサーボードをプログラムから動かす方法を学びます。温湿度気圧センサーの\ruby{情報}{じょう|ほう}を読み取ったり、センサーボードに取り付けてあるLEDを点灯・消灯させたり、ボタンを上手に使ったりして、センサーボードをプログラムから動かすことができるようになりましょう。

この章の後半には、コマンドをもっと使ってみたい人のための発展的な内容があります。
ファイルブラウザではできない、テキストの操作や、ファイルの生成などができるようになりましょう。

\end{document}
